% Sections won't have numbers
% Make the page area as large as possible
% Set the first level numbering to arabic and the second level
% to alphabetic. The remaining two levels are arabic
%Change to \pagestyle{empty} to suppress numbers on following pages
%\pagestyle{empty}
% Do not reset outermost enumerate counter


\documentclass[onecolumn]{article}
%%%%%%%%%%%%%%%%%%%%%%%%%%%%%%%%%%%%%%%%%%%%%%%%%%%%%%%%%%%%%%%%%%%%%%%%%%%%%%%%%%%%%%%%%%%%%%%%%%%%%%%%%%%%%%%%%%%%%%%%%%%%%%%%%%%%%%%%%%%%%%%%%%%%%%%%%%%%%%%%%%%%%%%%%%%%%%%%%%%%%%%%%%%%%%%%%%%%%%%%%%%%%%%%%%%%%%%%%%%%%%%%%%%%%%%%%%%%%%%%%%%%%%%%%%%%
\usepackage{amsmath}
\usepackage{sw20exm2}
\usepackage[compat2]{geometry}

\setcounter{MaxMatrixCols}{10}
%TCIDATA{OutputFilter=LATEX.DLL}
%TCIDATA{Version=5.50.0.2953}
%TCIDATA{<META NAME="SaveForMode" CONTENT="1">}
%TCIDATA{BibliographyScheme=Manual}
%TCIDATA{Created=Sunday, August 26, 2012 12:26:36}
%TCIDATA{LastRevised=Wednesday, July 01, 2026 14:53:34}
%TCIDATA{<META NAME="ViewSettings" CONTENT="0">}
%TCIDATA{<META NAME="GraphicsSave" CONTENT="32">}
%TCIDATA{<META NAME="DocumentShell" CONTENT="Exams and Syllabi\SW\SW Exam #1 - 8.5x14 Page">}
%TCIDATA{CSTFile=Exam.cst}

\setlength{\topmargin}{-0.75in}
\setlength{\textheight}{12.25in}
\setlength{\oddsidemargin}{0.0in}
\setlength{\evensidemargin}{0.0in}
\setlength{\textwidth}{6.5in}
\def\labelenumi{\arabic{enumi}.}
\def\theenumi{\arabic{enumi}}
\def\labelenumii{(\alph{enumii})}
\def\theenumii{\alph{enumii}}
\def\p@enumii{\theenumi.}
\def\labelenumiii{\arabic{enumiii}.}
\def\theenumiii{\arabic{enumiii}}
\def\p@enumiii{(\theenumi)(\theenumii)}
\def\labelenumiv{\arabic{enumiv}.}
\def\theenumiv{\arabic{enumiv}}
\def\p@enumiv{\p@enumiii.\theenumiii}
\pagestyle{plain}
\setcounter{secnumdepth}{0}
\newtheorem{theorem}{Theorem}
\newtheorem{acknowledgement}[theorem]{Acknowledgement}
\newtheorem{algorithm}[theorem]{Algorithm}
\newtheorem{axiom}[theorem]{Axiom}
\newtheorem{case}[theorem]{Case}
\newtheorem{claim}[theorem]{Claim}
\newtheorem{conclusion}[theorem]{Conclusion}
\newtheorem{condition}[theorem]{Condition}
\newtheorem{conjecture}[theorem]{Conjecture}
\newtheorem{corollary}[theorem]{Corollary}
\newtheorem{criterion}[theorem]{Criterion}
\newtheorem{definition}[theorem]{Definition}
\newtheorem{example}[theorem]{Example}
\newtheorem{exercise}[theorem]{Exercise}
\newtheorem{lemma}[theorem]{Lemma}
\newtheorem{notation}[theorem]{Notation}
\newtheorem{problem}[theorem]{Problem}
\newtheorem{proposition}[theorem]{Proposition}
\newtheorem{remark}[theorem]{Remark}
\newtheorem{solution}[theorem]{Solution}
\newtheorem{summary}[theorem]{Summary}
\newenvironment{proof}[1][Proof]{\noindent\textbf{#1.} }{\ \rule{0.5em}{0.5em}}
\input{tcilatex}
\begin{document}


\begin{center}
\begin{equation*}
\FRAME{itbpF}{1.0525in}{0.7515in}{0in}{}{}{Figure}{\special{language
"Scientific Word";type "GRAPHIC";maintain-aspect-ratio TRUE;display
"USEDEF";valid_file "T";width 1.0525in;height 0.7515in;depth
0in;original-width 2.3168in;original-height 1.6466in;cropleft "0";croptop
"1";cropright "1";cropbottom "0";tempfilename
'TCTQIQ04.bmp';tempfile-properties "XPR";}}
\end{equation*}%
\textbf{F\'{\i}sica II (2026-1)}

\textbf{Tarea 4}\emph{: Ley de Biot-Savart, Ley de Ampere y Ley de
Faraday-Lenz}

\textit{Profesor: Arturo Fern\'{a}ndez P\'{e}rez}

\textbf{Plazo de Entrega: Martes 07 de julio a las 20:00 hrs.}
\end{center}

\begin{enumerate}
\item Tres conductores rectos y muy extensos, conducen corrientes $I_{1}=2$
[A], $I_{2}=5$ [A] e $I_{3}=3$ [A], como muestra la Fig. 1. Considerando que 
$a=5$ [cm], $b=2$ [cm] y $c=7$ [cm] marcan las distancias al punto P y
utilizando el sistema de referencia indicado, determine:

\begin{enumerate}
\item El campo magn\'{e}tico total $\vec{B}$ en el punto P.

\item La fuerza magn\'{e}tica $\vec{F}_{m}$ ejercida sobre una secci\'{o}n
de $2$ [m] de longitud del lazo de corriente $I_{2}$, ejercida por la
corriente $I_{1}$.%
\begin{equation*}
\FRAME{itbpFU}{2.6662in}{2.6308in}{0in}{\Qcb{Figura 1}}{}{Figure}{\special%
{language "Scientific Word";type "GRAPHIC";maintain-aspect-ratio
TRUE;display "USEDEF";valid_file "T";width 2.6662in;height 2.6308in;depth
0in;original-width 8.3169in;original-height 8.2071in;cropleft "0";croptop
"1";cropright "1";cropbottom "0";tempfilename
'THI77900.bmp';tempfile-properties "XPR";}}
\end{equation*}
\end{enumerate}

\item Considere tres corrientes el\'{e}ctricas entrando o saliendo del plano
de la pantalla, como indica la Fig. 2. Si \'{e}stas se ubican en los v\'{e}%
rtices de un tri\'{a}ngulo rect\'{a}ngulo, determine el campo magn\'{e}tico
total $\vec{B}\,$en el centro de la hipotenusa, si $I_{1}=0,35$ [A], $%
I_{2}=2I_{1}$ e $I_{3}=3I_{1}$ (Indicaci\'{o}n: Utilize la ley de Ampere y
el vector unitario polar $\pm \phi $ para indicar la direcci\'{o}n del campo
magn\'{e}tico)%
\begin{equation*}
\FRAME{itbpFU}{2.3774in}{1.9052in}{0in}{\Qcb{Figura 2}}{}{Figure}{\special%
{language "Scientific Word";type "GRAPHIC";maintain-aspect-ratio
TRUE;display "USEDEF";valid_file "T";width 2.3774in;height 1.9052in;depth
0in;original-width 7.8784in;original-height 6.3045in;cropleft "0";croptop
"1";cropright "1";cropbottom "0";tempfilename
'THI7RK02.bmp';tempfile-properties "XPR";}}
\end{equation*}%
\pagebreak 

\item Los lazos de corriente de la Fig. 3 conducen ambos una corriente $I=1,5
$ [A] y sus arcos circulares subtienden un \'{a}ngulo de 30$%
%TCIMACRO{\U{b0}}%
%BeginExpansion
{{}^\circ}%
%EndExpansion
$, como muestra la Fig. 3. Si $a=5$ [cm], $b=3$ [cm]\thinspace\ y
considerando los sentidos de circulaci\'{o}n de la corriente indicados,
determine el campo magn\'{e}tico $\vec{B}$ en el punto $O.$ (Indicaci\'{o}n:
Utilize la forma escalar de la ley de Biot-Savart y la regla de la mano
derecha)%
\begin{equation*}
\FRAME{itbpFU}{2.6991in}{1.6172in}{0in}{\Qcb{Figura 3}}{}{Figure}{\special%
{language "Scientific Word";type "GRAPHIC";maintain-aspect-ratio
TRUE;display "USEDEF";valid_file "T";width 2.6991in;height 1.6172in;depth
0in;original-width 6.5726in;original-height 3.9271in;cropleft "0";croptop
"1";cropright "1";cropbottom "0";tempfilename
'THI8LU03.bmp';tempfile-properties "XPR";}}
\end{equation*}

\item Una espira rect\'{a}ngular de \'{a}rea $A=50$ [cm$^{2}$] gira con una
velocidad angular $\omega =720$ [$rpm$] en una zona donde hay un campo magn%
\'{e}tico uniforme de magnitud $B=$ $0.5$ $[T]$, como muestra la Fig. 4.
Determine la fem inducida $\varepsilon $ en la espira.%
\begin{equation*}
\FRAME{itbpFU}{1.8178in}{1.9415in}{0in}{\Qcb{Figura 4}}{}{Figure}{\special%
{language "Scientific Word";type "GRAPHIC";maintain-aspect-ratio
TRUE;display "USEDEF";valid_file "T";width 1.8178in;height 1.9415in;depth
0in;original-width 5.6464in;original-height 6.0364in;cropleft "0";croptop
"1";cropright "1";cropbottom "0";tempfilename
'THI95804.bmp';tempfile-properties "XPR";}}
\end{equation*}

\item Una espira de $N=20$ vueltas, radio $R=8$ [cm] y resistencia el\'{e}%
ctrica de 20 [$\Omega $] se ubica en el centro de un solenoide muy largo, el
cu\'{a}l tiene $n=200$ [vueltas/metro], radio $r=5$ [cm] y es conc\'{e}%
ntrico a la espira, como se ve en la Fig. 5. Si por el solenoide circula una
corriente $I(t)=\frac{5}{\sqrt{t}}$ [A], determine,

\begin{enumerate}
\item El flujo magn\'{e}tico $\phi _{B}$ a trav\'{e}s de la espira

\item La magnitud de la fem $\varepsilon $ y la corriente inducida $i(t)$ en
la espira \textquestiondown La corriente inducida $i(t)$ circula en la misma
direcci\'{o}n que la corriente del solenoide $I(t)$? Justifique.%
\begin{equation*}
\FRAME{itbpFU}{1.3041in}{1.6942in}{0in}{\Qcb{Figura 5}}{}{Figure}{\special%
{language "Scientific Word";type "GRAPHIC";maintain-aspect-ratio
TRUE;display "USEDEF";valid_file "T";width 1.3041in;height 1.6942in;depth
0in;original-width 2.2295in;original-height 2.9066in;cropleft "0";croptop
"1";cropright "1";cropbottom "0";tempfilename
'THI7EE01.wmf';tempfile-properties "XPR";}}
\end{equation*}
\end{enumerate}
\end{enumerate}

\end{document}
